% Exemplo: definir comandos próprios com \newcommand (sem e com argumentos)
% Compilar: latexmk -pdf comandos-proprios.tex     Testado com: TeX Live 2024 ou mais recente
\documentclass[a4paper]{article}
\usepackage[T1]{fontenc}
\usepackage{lmodern}
\usepackage[brazilian]{babel}
%%BEGIN:comandospreambulo
% sem argumentos: um nome que se repete muito no texto
\newcommand{\sistema}{\textsf{GestaFácil}}
% com um argumento (#1): termos em outro idioma sempre em itálico
\newcommand{\estrangeiro}[1]{\emph{#1}}
% com um argumento opcional (padrão "pendente") e um obrigatório
\newcommand{\revisar}[2][pendente]{\textbf{[#1: #2]}}
%%END:comandospreambulo

\begin{document}

%%BEGIN:comandosuso
O \sistema{} foi desenvolvido com o \estrangeiro{framework} Django. A versão
móvel do \sistema{} ainda está em teste \revisar{conferir a data}. A interface
usa \estrangeiro{layout} responsivo \revisar[ok]{testado em 3 navegadores}.
%%END:comandosuso

\end{document}
